\section{The cost of exact balance}
\label{sec:theory}
In this section $k=n$. A block is \emph{position-balanced} if each option appears in each slot
once. Writing its orders as rows gives a Latin square. For a single item
with nonconstant profile $d$, divide the variance of the block mean
by the population variance $v(d)$ across all $n!$ orders. Denote the
ratio by $R_D(d)$ for design $D$. A ratio of one is the variance of
evaluating one uniform order. The minimax risk is the smallest attainable
supremum of $R_D(d)$ over the indicated class of profiles. A correctness
difference can take any values in $\{-1,0,1\}$ by assigning the two
models suitable binary outcomes at each order.

\paragraph{Theorem 1 (position balance).}
For every $n\ge4$, every position-balanced design with uniform order
marginals has a correctness difference, depending on the displayed
positions of just two options, whose normalized risk is at least
\begin{equation}
 L_n=\frac{n-2}{n^2-3n+1}.
 \label{eq:two-option}
\end{equation}
This exceeds the unrestricted optimum $(n!-n)/(n(n!-1))$ for the same
budget. For $6\le n\le10$ and for $n=2p$ with any prime $p\ge131$ there is a design whose risk is at most
$L_n$ for \emph{every real profile}. Hence $L_n$ is the exact minimax
risk at all these sizes for either contrast class.

\emph{Proof idea.} Fix two semantic options. Their ordered pair of slots
has $m=n(n-1)$ possible values. A balanced block visits $n$ of them,
with each slot used once by each option. Subtract the uniform dyad
weights from the empirical weights and call the result $e$.
Its row and column sums vanish, leaving $d_0=n^2-3n+1$ dimensions.
Every block has $m\|e\|^2=n-2$. Thus the covariance on this space has
a fixed trace, and its largest eigenvalue is at least its average
eigenvalue $L_n$ after normalization. For correctness profiles,
an alternating signed four-cycle on slots and an oriented triangle
give two ternary contrasts. Relabeling their slots and mixing them
in proportion to the dimensions of the symmetric and skew-symmetric
subspaces makes their average normalized risk exactly $L_n$.
At least one realizes the lower bound or more. The upper certificates
use two Latin squares at each $n$ and bound the full permutation
spectrum. For the infinite family, the two squares have just three relative-row cycle types. Mixing them equalizes the two dyad risks, and a uniform character bound controls every other direction (Appendix~\ref{app:two-prime}). The small-size certificates are in Appendix~\ref{app:general-position}.

The two-option argument extends to any position-balanced budget $k$.
Set $m=n(n-1)$ and let $h$ be the remainder of $k$ upon division by
$m$. Some correctness profile has risk at least
$h(m-h)/(k^2 d_0)$. This recovers Eq.~\eqref{eq:two-option} at $k=n$
and can be zero at budgets divisible by $m$
(Appendix~\ref{app:general-position}).

\paragraph{A strict-cost test for correctness.}
Take two distinct rows uniformly from a sampled block and let $t$ be the
expected number of two-cycles in their relative permutation. A two-cycle
means that two options have swapped slots. Write $s=n(n-3)/2$,
$a=(n-1)(n-2)/2$. The same two ternary witnesses give
\begin{equation}
 \sup_d R_D(d)\ \ge\ \frac1n+\frac{n-1}{n}
       \max\!\left\{\frac{t}{s},\frac{1-t}{a}\right\}.
 \label{eq:swaps}
\end{equation}
The right side equals $L_n$ only at $t=s/(s+a)$. If all blocks
admitted by an extra rule have swap averages strictly on one side
of this target, and position balance attains $L_n$, the rule
strictly increases correctness minimax risk. This test needs only
pairs of rows. At four options the two Latin types have $t=2/3$
and $2$, giving the sharp position lower bound $1/2$. At six,
every row-complete block has $t=3/5$, already certifying a strict
increase over $4/19$ before computing the exact cost below.

For real scores the same proof tests additional balance requirements.
If every block integrates $r$ extra functions of these two slots that
are linearly independent modulo row and column effects, its covariance
has at most $d_0-r$ remaining dimensions. Its normalized worst-case
risk is at least $(n-2)/(d_0-r)$ when $r<d_0$.
Whenever position balance attains $L_n$, every such restriction with
$r>0$ has a strictly positive cost. The direction of the effect on a
particular pair of models depends on its response profile.

\begin{table}[t]
\centering
\caption{Normalized minimax variance for differences of correctness.
Each model gets $n$ calls. The last column requires every directed
adjacency once per block, except at five options, where 15 of 20
is the maximum. Brackets give certified bounds.}
\label{tab:all-n}
\small
\begin{tabular}{rcc}
\toprule
Options & Position balance & With adjacency \\
\midrule
4 & $1/2$ & $1$ \\
5 & $9/25$ & $[7/15,\;0.47087]$ \\
6 & $4/19$ & $1$ \\
7 & $5/29$ & full coverage infeasible \\
8 & $6/41$ & $[69/280,\;1)$ \\
9 & $7/55$ & $[43/243,\;1)$ \\
10 & $8/71$ & $[1/8,\;1)$ \\
\bottomrule
\end{tabular}
\end{table}

\subsection{When full adjacency costs all variance reduction}
\label{sec:connectivity}
Call a position-balanced block \emph{row-complete} if every ordered
pair of distinct options is consecutive in some row. There are
$n(n-1)$ adjacencies available across $n$ rows, so each must occur
exactly once. Row-complete Latin squares are classical
\citep{williams1949,wanlesscatalogue}. We use the orders that
can be sampled together to decide their risk.

\paragraph{Theorem 2 (connectivity).}
Let an allowed family of $k$-order blocks on $S_n$ be nonempty and
closed under relabeling the options. Connect two orders if they
occur together in an allowed block. The minimax risk equals one
if and only if this graph is disconnected, for real profiles and
for correctness differences. If it is disconnected, every feasible
unbiased design has worst-case risk one. The components are the left
cosets of the subgroup generated by rows of allowed blocks normalized
to contain the identity (Appendix~\ref{app:connectivity}).

\emph{Proof idea.} Draw a block and then a uniform row in it. The row
is uniform over all orders. The variance of its score is the variance of
the block mean plus the expected within-block variance. On a disconnected
graph, the indicator of a component is constant on every allowed
block. All its variance survives. On a connected graph, the uniform
law on allowed blocks is unbiased, and every nonconstant profile varies
within some block. Compactness on the centered unit sphere gives a
strict bound below one.

With six options, exactly two of the 9,408 reduced Latin squares are
row-complete. Their rows generate a subgroup of size 24,
splitting the 720 orders into 30 components. For a concrete witness,
take the indicator that two designated options occupy opposite slots
(first and last, second and fifth, or third and fourth). Every row of
a row-complete block gives the same answer. This indicator has mean
$1/5$, and all of its variance survives the six calls. Requiring adjacency raises
the optimum from $4/19$ to $1$, a factor of $4.75$.
At four options there are six components. At eight, two valid squares
already generate all of $S_8$. At nine, one suffices. Connectivity
allows variance reduction in these larger designs, although the
worst-case cost of adjacency remains strictly positive
(Table~\ref{tab:all-n}).

\paragraph{Quantifying a partial loss.}
For a set $A$ of orders, put $p=|A|/n!$ and $s_B=|A\cap B|$.
The same variance identity gives the exact risk of its indicator,
\begin{equation}
 R_D(\mathbf1_A)=1-
 \frac{\mathbb E_D[s_B(k-s_B)]}{k^2 p(1-p)},\qquad 0<p<1 .
 \label{eq:cut}
\end{equation}
More generally, average the normalized risks of several feasible
contrasts. If this average exceeds a reference risk on every
option-relabeling orbit of eligible blocks, every design has a
worst-case risk above the reference. Relabeling an arbitrary design
can only decrease its worst-case risk. This gives a quantitative
test even when the co-occurrence graph is connected.

At five options position balance alone has sharp risk $9/25$ for
both contrast classes. The permutation-sign contrast has block mean of
magnitude at least $3/5$ in every Latin block, and one doubly relabeled
template attains the lower bound. At maximum coverage, three ternary
contrasts with weights $1/45$, $4/9$,
and $8/15$ give the same weighted risk $7/15$ on \emph{each} of the
five maximum-coverage orbits (Appendix~\ref{app:five-certificates}).
Thus the relative cost over $9/25$ is at least $8/27$ (about $29.6\%$).
The certified upper bound $0.47087$ leaves less than $0.00421$
between bounds. For unrestricted real scores the interval is
$[0.47084,0.47087]$. The exact optimum for correctness at five remains
open. At eight through ten, witnesses depending on only two options
already separate row completeness from the position optimum.
Appendix~\ref{app:connectivity} gives their construction and the
catalogue inputs.
